%!TEX root = ./main.tex

\chapter{Metodología y Diseño} \label{ch:metodologia}
\graphicspath{{figs/}}

% Breve introducción del capítulo indicando qué se va a explicar (el flujo de trabajo y las decisiones arquitectónicas).
El presente capítulo detalla la metodología de diseño e implementación del sistema HR-RT-Reporter...

% =======================================================================
% SECCIÓN 1: Arquitectura General
% =======================================================================
\section{Arquitectura General del Sistema}
% Explicar la visión global: el SoC Zynq-7000, la placa Zybo, y por qué se divide en dos mundos (PS y PL).

\subsection{Separación de Dominios: Sistema de Procesamiento (PS) y Lógica Programable (PL)}
\label{ssec:separacion_dominios}

Para cumplir simultáneamente con los requisitos de determinismo temporal del software crítico y de trazabilidad no intrusiva de eventos, la arquitectura del sistema propuesto implementa una estricta separación de responsabilidades físicas y lógicas. Esta división aprovecha de manera óptima las dos divisiones principales del SoC Zynq-7000: el Sistema de Procesamiento (PS, \textit{Processing System}) y la Lógica Programable (PL, \textit{Programmable Logic}).

En el dominio del \textbf{Sistema de Procesamiento (PS)}, el microprocesador ARM Cortex-A9 ejecuta una distribución del sistema operativo Linux empotrado. Este procesador actúa como el elemento maestro y supervisor de la plataforma experimental, asumiendo las siguientes responsabilidades clave:
\begin{itemize}
    \item \textbf{Control de la Ejecución:} Coordina y habilita la inicialización o detención del procesador \textit{soft-core} en la PL mediante la gestión del estado lógico de su señal de reinicio física (\textit{reset}), controlada a través de un canal AXI GPIO de un bit.
    \item \textbf{Recolección Asíncrona de Datos:} Lee de manera periódica y diferida los registros de marcas temporales almacenados por la lógica del hardware de medición.
    \item \textbf{Gestión y Exportación:} Procesa, organiza y exporta la información capturada hacia el espacio de usuario (\textit{user-space}) en el almacenamiento de Linux para su análisis y evaluación fuera de línea, eliminando la necesidad de realizar procesamiento intrusivo de trazabilidad en tiempo real.
\end{itemize}

Por su parte, en el dominio de la \textbf{Lógica Programable (PL)} residen los componentes de hardware que garantizan la predictibilidad matemática del software crítico bajo ensayo:
\begin{itemize}
    \item \textbf{Procesador Soft-core (Ibex):} Núcleo basado en la arquitectura RISC-V que ejecuta de manera exclusiva la aplicación bajo prueba. Al ejecutarse directamente en el tejido físico de la FPGA de forma aislada, es inmune a las interrupciones o cambios de contexto impredecibles característicos del sistema operativo del procesador ARM.
    \item \textbf{Memoria Local (BRAM):} Implementada mediante bloques de RAM embebida dentro del propio integrado Zynq. Está configurada en modo \textit{True Dual Port} para satisfacer las exigencias de una arquitectura tipo Harvard, permitiendo que el procesador Ibex acceda concurrentemente al flujo de instrucciones de arranque y de datos de control en un único ciclo de reloj ($20\text{ ns}$ a $50\text{ MHz}$).
    \item \textbf{Módulo de Monitoreo (Reporter):} Es el bloque personalizado que implementa el contador de alta velocidad y la lógica necesaria para espiar el bus de datos del SUT de forma transparente. Al detectar operaciones de almacenamiento específicas, captura los valores temporales exactos del reloj de hardware sin interponer demoras en el código bajo ejecución.
\end{itemize}

La comunicación e intercambio de información entre ambos mundos (\textit{hard-core} en el PS y \textit{soft-core} en la PL) se canaliza mediante el estándar de interconexión \textbf{AMBA AXI4} (principalmente en su variante optimizada para registros AXI4-Lite). A través de esta red, el procesador ARM mapea las direcciones físicas de la memoria local BRAM y los registros de estado del Reporter dentro de su propio mapa de direccionamiento virtual. Esto permite transferir comandos de arranque y descargar registros con alta velocidad, de forma controlada y sin degradar la fidelidad del análisis temporal.
% =======================================================================
% SECCIÓN 2: Preparación e Integración del Núcleo RISC-V
% =======================================================================
\section{Diseño e Integración del Subsistema RISC-V}

La implementación física del Sistema Bajo Prueba (SUT) requiere instanciar el procesador \textit{soft-core} dentro del tejido lógico de la FPGA. Sin embargo, antes de proceder a la síntesis en Vivado, el código fuente del núcleo debe ser correctamente estructurado y empaquetado. 

\subsection{Entorno de Desarrollo y Empaquetado del Núcleo}

La preparación e integración del procesador Ibex requirió el establecimiento de un entorno de desarrollo basado en Linux. Para tal fin, se empleó el Subsistema de Windows para Linux (WSL, por sus siglas en inglés), lo que permitió operar con herramientas nativas de este sistema operativo sin necesidad de abandonar el sistema anfitrión.

El núcleo Ibex, mantenido bajo el estándar de OpenTitan, consta de una vasta jerarquía de archivos fuente escritos en SystemVerilog. Debido a la complejidad arquitectónica y a la gran cantidad de dependencias cruzadas e inclusiones (\textit{includes}), la importación manual de estos componentes a un entorno de síntesis resulta altamente ineficiente y propensa a errores. Para resolver esta problemática, se implementó el uso de \textbf{FuseSoC}, un gestor de paquetes y herramienta de construcción (\textit{build tool}) especializado en lenguajes de descripción de hardware (HDL).

%A fines de garantizar la estabilidad, portabilidad y reproducibilidad del entorno de desarrollo, la instalación de la herramienta FuseSoC se llevó a cabo dentro de un entorno virtual de Python. Esta práctica de aislamiento previene conflictos entre las bibliotecas requeridas por el compilador HDL y los paquetes globales del sistema operativo anfitrión.

\subsection{FuseSoC y Gestión de Dependencias para Vivado}
\label{ssec:fuse}
FuseSoC es una herramienta de código abierto diseñada para simplificar la gestión, reutilización e integración de núcleos de Propiedad Intelectual (IP) en proyectos de diseño digital. Su función principal es proveer un sistema unificado de administración de dependencias y automatización de compilación (\textit{build system}), operando de manera independiente del fabricante y de la herramienta EDA utilizada \cite{kindgren2019scalable}.

La herramienta organiza cada núcleo mediante un archivo de descripción de núcleo (\textit{core description file}), en el cual se especifican de forma explícita los archivos fuente, las dependencias internas, los parámetros de configuración y los objetivos de compilación o simulación. A partir de esta especificación, FuseSoC resuelve de manera determinística las relaciones entre módulos y genera la estructura de archivos requerida por herramientas de síntesis.

Una vez configurado el entorno, se procedió a descargar el código fuente desde el repositorio oficial de Ibex. Posteriormente, se ejecutó FuseSoC apuntando al directorio raíz del proyecto. La herramienta analizó el manifiesto del núcleo, resolviendo de manera automatizada la jerarquía de los archivos, el orden lógico de compilación y todas las dependencias estructurales requeridas. 

El resultado de este proceso de empaquetado fue la generación de un directorio de salida limpio (\texttt{build/ibex\_out/src/}), el cual consolidó exclusivamente los archivos SystemVerilog necesarios para el diseño. Dentro de este conjunto de archivos preprocesados, se expuso a \texttt{ibex\_top.sv} como el archivo de cabecera o entidad de nivel superior (\textit{top-level module}). Finalmente, este directorio fue exportado e importado como fuente de diseño en el entorno de AMD Vivado, permitiendo instanciar el núcleo RISC-V como un bloque funcional integro.

\subsection{Encapsulamiento del Núcleo (Wrapper)}
\label{ssec:wrapper_ibex}

El núcleo RISC-V original (\texttt{ibex\_top.sv}), empaquetado y exportado mediante la herramienta FuseSoC, posee una extensa cantidad de puertos de entrada y salida (I/O). Gran parte de estos pines están diseñados para extensiones del estándar, depuración o verificación formal (RVFI), características que no son requeridas en la arquitectura de este proyecto. Dejar dichos pines desconectados o exponerlos directamente en el entorno de diseño por bloques de Vivado (\textit{Block Design}) resulta en una fuente recurrente de errores de síntesis (\textit{unconnected pins errors}) y genera un esquema visualmente inmanejable.

Para mitigar esta problemática y garantizar una integración limpia en el hardware físico, se desarrolló una estrategia de encapsulamiento jerárquico estructurada en dos niveles modulares:

\begin{itemize}
    \item \textbf{Nivel de Subsistema (\texttt{ibex\_bram\_subsystem.sv}):} Este módulo de hardware, escrito en SystemVerilog, actúa como el primer nivel de abstracción. Su función principal es alojar internamente al procesador Ibex puro y acoplarlo físicamente con su memoria local (BRAM). En este bloque se enmascaran y terminan lógicamente todos los puertos I/O no utilizados del núcleo, resolviendo la complejidad interna de las conexiones. Como resultado, se consolida un subsistema de ``procesador más memoria'' completamente autocontenido.
    
    \item \textbf{Nivel de Integración (\texttt{ibex\_bram\_wrapper\_bd.v}):} Debido a que la herramienta \textit{IP Integrator} de Vivado presenta un comportamiento de síntesis y ruteo más robusto al interactuar con módulos de nivel superior (\textit{top-level}) descritos en lenguaje Verilog estándar, se desarrolló este segundo envoltorio (\textit{wrapper}). Este archivo instancia exclusivamente a \texttt{ibex\_bram\_subsystem.sv} y tiene el propósito estructural de exponer hacia el entorno gráfico del \textit{Block Design} únicamente las señales esenciales: los relojes, las señales de reinicio (\texttt{reset}) y las interfaces de datos dedicadas a la comunicación con el módulo Reporter.
\end{itemize}

Esta arquitectura jerárquica simplifica drásticamente la integración funcional del procesador dentro de la Lógica Programable (PL), previniendo errores de compilación y asegurando que las herramientas de AMD Xilinx solo gestionen las rutas de interconexión estrictamente necesarias para el ensayo.

\subsection{Acondicionamiento de Reloj (Clock Wizard)}
\label{ssec:clock_wizard}

La estabilidad y precisión de la señal de sincronismo constituyen un requisito crítico y mandatorio para la plataforma HR-RT-Reporter. Dado que el sistema de trazabilidad calcula el desempeño del software mediante un contador de hardware de alta resolución, cualquier fluctuación frecuencial o fenómeno de \textit{jitter} degradaría la validez matemática de las marcas temporales (\textit{timestamps}) obtenidas. Para subsanar esta vulnerabilidad física, se integró el bloque de Propiedad Intelectual (IP) de AMD Xilinx denominado \textbf{Clocking Wizard}.

El \textit{Clocking Wizard} utiliza los recursos lógicos analógicos dedicados del SoC Zynq-7000, específicamente un bucle de enganche de fase (PLL, \textit{Phase-Locked Loop}) y un administrador de reloj en modo mixto (MMCM, \textit{Mixed-Mode Clock Manager}). Este bloque toma como referencia el oscilador de cristal externo de $125\text{ MHz}$ de la placa de desarrollo Zybo y sintetiza una señal de reloj de salida estable y dedicada de $50\text{ MHz}$ para el dominio de la Lógica Programable (PL).

La elección de una frecuencia de operación de $50\text{ MHz}$ obedece a dos factores de diseño fundamentales:
\begin{enumerate}
    \item \textbf{Resolución Temporal Estricta:} Permite que el contador interno del Módulo Reporter se incremente exactamente cada un ciclo de reloj, equivalente a un período discreto de $20\text{ ns}$. Esta resolución es más que suficiente para caracterizar las transiciones de estado en algoritmos de control de tiempo real.
    \item \textbf{Relajación de Restricciones Temporales (\textit{Timing Constraints}):} Al operar a una frecuencia moderada en comparación con los límites del silicio, se facilita el proceso de enrutamiento y síntesis física de Vivado. Esto reduce drásticamente el riesgo de violación de tiempos de establecimiento (\textit{setup}) y retención (\textit{hold}) en los caminos críticos entre el procesador Ibex y el bus AXI.
\end{enumerate}

Un aspecto metodológico clave de esta configuración es la gestión de la señal de retroalimentación \texttt{locked} provista por el IP. Al energizar el sistema o al reconfigurar la FPGA, las salidas de reloj del MMCM experimentan un régimen transitorio inestable en el que la frecuencia de salida oscila antes de estabilizarse. La señal de hardware \texttt{locked} permanece en bajo (`0`) durante esta transición y se afirma en alto (`1`) de forma asíncrona únicamente cuando la señal de salida cumple con los márgenes estrictos de fase y frecuencia configurados.

Para evitar que el subsistema de procesamiento inicie operaciones en un estado de reloj indeterminado (lo cual provocaría comportamientos erráticos o bloqueos lógicos irreversibles en el procesador RISC-V), se diseñó una barrera de seguridad lógica en el código de inicialización del sistema. La señal de reset interno del subsistema (\texttt{rst\_n}) se condiciona mediante una compuerta evaluada en hardware que combina la señal de reset controlada por el ARM (\texttt{rst\_pin}) y la señal de estabilidad del reloj (\texttt{locked}), como se describe en la ecuación \ref{eq:reset_logic}:

\begin{equation}
    \label{eq:reset_logic}
    \text{rst\_n} = \neg(\text{rst\_pin}) \wedge \text{locked}
\end{equation}
                                                                            CHEQUEAR!!!!!!!!!!
                                                                            
Esta implementación garantiza que el procesador Ibex permanezca físicamente retenido en estado de reinicio activo bajo hasta que la señal de $50\text{ MHz}$ se encuentre completamente estabilizada y libre de transitorios. De este modo, la ejecución del firmware crítico inicia bajo condiciones de sincronismo óptimas y predecibles.

% =======================================================================
% SECCIÓN 3: Arquitectura de Memoria y Arbitraje
% =======================================================================
\section{Arquitectura de Memoria y Prevención de Colisiones}
% Explicar cómo se le da memoria al Ibex y cómo se protege de accesos indebidos.
\subsection{Memoria BRAM en Modo True Dual Port}
\label{ssec:bram_dual_port}

Para satisfacer las exigencias de baja latencia y aislar al núcleo Ibex de la congestión del bus principal del sistema (DDR3 externa), se optó por implementar la memoria local utilizando la primitiva \textit{Block RAM} (BRAM) interna de la FPGA. Esta memoria estática permite accesos en un único ciclo de reloj ($20\text{ ns}$ a $50\text{ MHz}$), asegurando un determinismo absoluto en la ejecución algorítmica.

El procesador Ibex está basado en una arquitectura tipo Harvard, por lo que requiere buses físicamente separados para el flujo de instrucciones (\textit{Instruction Fetch}) y el flujo de datos (\textit{Load/Store}). Para acomodar esta topología en un único bloque físico, la BRAM fue configurada en modo \textbf{True Dual Port} (Doble Puerto Verdadero). La asignación lógica de los puertos se definió de la siguiente manera:
\begin{itemize}
    \item \textbf{Puerto A (Solo Lectura para SUT):} Se dedica exclusivamente a la obtención de instrucciones.
    \item \textbf{Puerto B (Lectura/Escritura para SUT):} Se utiliza para la carga y almacenamiento de variables o punteros generados durante la ejecución del programa.
\end{itemize}

Además, la herramienta de síntesis fue configurada para aceptar un archivo de inicialización de coeficientes (\texttt{.coe}), el cual permite que la memoria BRAM se instancie con el firmware del Ibex pre-cargado desde el momento exacto en que se programa la FPGA (\textit{bitstream}), eliminando la necesidad de que el ARM transfiera el código de arranque rutinariamente.

\subsubsection{Dimensionamiento y Configuración de Celdas}
\label{sssec:bram_dimensions}

El bloque de memoria local del procesador de prueba fue dimensionado de manera estricta para optimizar el uso de los recursos lógicos de la FPGA sin comprometer el espacio necesario para el firmware de ensayo. Las especificaciones de configuración del silicio se detallan a continuación:

\begin{itemize}
    \item \textbf{Ancho de Palabra del Bus (\textit{Data Width}):} Configurado nativamente en \textbf{$32\text{ bits}$}. Esto coincide exactamente con el tamaño de palabra de la arquitectura de la ISA RISC-V (RV32I) implementada por el núcleo Ibex, permitiendo transferencias completas de instrucciones y datos en un único ciclo sin necesidad de deserialización.
    \item \textbf{Profundidad de Celdas (\textit{Memory Depth}):} Se fijó una profundidad de \textbf{$2048$ celdas} (posiciones de memoria discretas).                                                                                             %<<<<-------------------- REVISAR LOS DATOS.
    \item \textbf{Capacidad Total de Almacenamiento:} El producto de las dimensiones físicas ($2048 \text{ posiciones} \times 32 \text{ bits}$) equivale a un total de $65,536\text{ bits}$ de memoria interna. En términos de almacenamiento direccionable por bytes, esto representa una capacidad neta de \textbf{$8\text{ KiB}$}, dimensionada en estricta consonancia con las restricciones de direccionamiento asignadas en el bus del sistema.
\end{itemize}

\subsubsection{Mapeo Cruzado de Direcciones del Sistema}
\label{sssec:memory_map}

Debido a la naturaleza heterogénea de la plataforma, el mapa de memoria debe analizarse desde dos perspectivas lógicas completamente disjuntas. Por un lado, el procesador ARM (PS) visualiza el hardware embebido como periféricos esclavos mapeados en su espacio físico general a través del bus de interconexión AXI. Por otro lado, el núcleo Ibex (PL) posee su propio direccionamiento local de ejecución interna de $32\text{ bits}$ protegido por el hardware de exclusión mutua.

En la Tabla \ref{tab:mapa_memoria} se detallan las direcciones base, los tamaños lógicos asignados en el \textit{Address Editor} de AMD Vivado y las funciones de cada bloque desde la perspectiva de ambos dominios de procesamiento.

\begin{table}[htbp]
\centering
\caption{Mapa de direccionamiento de memoria unificado de la arquitectura heterogénea.}
\label{tab:mapa_memoria}
\begin{tabular}{llcl}
\hline
\textbf{Periférico / Interfaz} & \textbf{Master base Address}& \textbf{Range}& \textbf{Tamaño}\\ \hline
AXI BRAM Controller            & \texttt{0x40000000}                                                                 & $8\text{ KB}$   & \texttt{0x00000000} (Local BRAM)                                                        \\
AXI GPIO 0 (Reset SUT)         & \texttt{0x41200000}                                                                 & $64\text{ KB}$  & Inaccesible (Línea física \texttt{rst\_ni})                                            \\
AXI GPIO 1 (LEDs Debug)        & \texttt{0x41210000}                                                                 & $64\text{ KB}$  & Inaccesible (Línea de diagnóstico)                                                     \\
HR-RT-Reporter (IP Core)       & \texttt{0x43C00000}                                                                 & $64\text{ KB}$  & \texttt{0x80000000} (Registro Pasivo)                                                   \\ \hline
\end{tabular}
\end{table}

Como se observa en la distribución de la tabla, la ventana del controlador de memoria \textit{AXI BRAM Controller} está dimensionada con un rango lógico de asignación de $8\text{ KB}$ en el bus ARM (restringido desde la dirección \texttt{0x40000000} hasta \texttt{0x40002000}), lo cual se acopla de manera simétrica con la capacidad física real de la memoria local direccionada por el Ibex a partir de su dirección interna \texttt{0x00000000}.

De igual manera, el módulo de trazabilidad por hardware es accedido por el ARM para la extracción asíncrona de marcas de tiempo en la dirección física \texttt{0x43C00000}, mientras que el Ibex únicamente requiere escribir en la dirección de frontera \texttt{0x80000000} para activar instantáneamente los disparadores de captura sin provocar retrasos por contención de buses.

% =======================================================================
% COMPLEMENTO DE LA SECCIÓN 3: Exclusión Mutua
% =======================================================================
\subsection{Lógica de Exclusión Mutua (El MUX de Protección)}
\label{ssec:exclusion_mutua}

La coexistencia de dos procesadores con capacidades de acceso sobre el mismo bloque físico de memoria de doble puerto introduce un problema crítico de concurrencia y consistencia de datos. Si el sistema operativo Linux, ejecutado en el entorno ARM, intentara realizar una operación de lectura o escritura invasiva sobre la BRAM mientras el procesador Ibex se encuentra en pleno cómputo de un algoritmo de tiempo real, se producirían colisiones en los buses internos y condiciones de carrera (\textit{race conditions}). Este fenómeno corrompería de forma inmediata los registros del programa bajo análisis, invalidando la fidelidad del ensayo temporal.
Para erradicar de forma absoluta esta vulnerabilidad sin introducir penalizaciones por software (\textit{overhead} de ejecución derivado del uso de semáforos lógicos), se diseñó e implementó un \textbf{Multiplexor de Exclusión Mutua por Hardware}. Este circuito de arbitraje asíncrono reside en el nivel superior del subsistema encapsulado (\texttt{ibex\_bram\_subsystem.sv}) y actúa como un cortafuegos en el silicio.

El principio operativo del multiplexor se fundamenta en la monitorización continua del estado físico del procesador RISC-V a través de su señal de reinicio activo en bajo (\texttt{rst\_ni}). La lógica de control implementada discrimina el acceso a las líneas de habilitación de escritura del Puerto B (\texttt{bram\_we}) mediante la evaluación de las señales lógicas descritas en la ecuación \ref{eq:ternario_mux}:

\begin{equation}
    \label{eq:ternario_mux}
    \text{bram\_we} = (\text{data\_we} \wedge \neg\text{sel\_ts}) \text{ ? } \text{4'b1111} : \text{4'b0000}
\end{equation}

A través de esta expresión ternaria evaluada en hardware puro, el sistema garantiza un comportamiento determinístico dividido en dos fases operativas estrictas:
\begin{enumerate}
    \item \textbf{Fase de Inyección y Diagnóstico ($\text{rst\_ni} = 0$):} Cuando el procesador ARM retiene al núcleo Ibex en estado de reset, el multiplexor concede privilegios de lectura y escritura totales al bus maestro AXI controlado por Linux. Esto permite que el sistema operativo inicialice la memoria, verifique la integridad de las líneas funcionales mediante patrones conocidos (e.g., comprobación del vector \texttt{0xCAFEBABE}) y prepare el entorno de prueba.
    \item \textbf{Fase de Aislamiento y Ejecución ($\text{rst\_ni} = 1$):} En el instante exacto en que Linux libera el reset del soft-core para dar inicio al experimento, el multiplexor conmuta de forma asíncrona. El procesador ARM pierde instantáneamente todo privilegio de escritura sobre las celdas de la BRAM. Cualquier intento de acceso fortuito o erróneo desde el espacio de usuario (\textit{user-space}) en Linux es interceptado y lógicamente descartado por las compuertas lógicas, aislando por completo al bus local del Ibex (\texttt{data\_*}) y protegiendo el determinismo de la prueba con un impacto de intrusión nulo (efecto sonda cero).
\end{enumerate}

% =======================================================================
% SECCIÓN 4: Módulo de Trazabilidad (HR-RT-Reporter)
% =======================================================================
\section{Diseño del Módulo AXI Reporter}
\label{sec:axi_reporter}

El núcleo de la propuesta metodológica de este Proyecto Final Integrador recae en el desarrollo del módulo IP personalizado denominado \texttt{axi\_reporter}. Este componente de hardware tiene como única finalidad registrar, documentar y fechar eventos de ejecución de software crítico de manera completamente determinística y sin alterar el comportamiento temporal de la aplicación en el procesador RISC-V.

\subsection{Generación de Timestamps por Hardware}
\label{ssec:timestamps_hw}

La medición de tiempos en sistemas operativos de propósito general (como Linux ejecutándose en el procesador ARM) está sujeta a latencias variables causadas por la planificación de procesos (\textit{scheduling}), manejo de interrupciones y políticas de caché. Para erradicar estas imprecisiones, el módulo \texttt{axi\_reporter} implementa un cronómetro de hardware dedicado basado en un contador ascendente continuo de **$64\text{ bits}$**.

Este contador es alimentado directamente por la señal de reloj purificada de $50\text{ MHz}$ proveniente del bloque \textit{Clocking Wizard}, operando en estricto sincronismo con el núcleo Ibex. Como consecuencia, el registro temporal se incrementa de forma inexorable cada ciclo de reloj, proveyendo una base de tiempo monótona con una resolución discreta y absoluta de **$20\text{ ns}$**.

El contador de $64\text{ bits}$ se inicializa en cero durante el estado de reinicio general del sistema y, una vez liberado, es capaz de registrar un tiempo de ejecución continuo superior a $11,000$ años antes de sufrir un desbordamiento de entero (\textit{overflow}), lo que garantiza una cobertura temporal más que suficiente para cualquier ensayo de verificación en software de misión crítica.

\subsection{Mecanismo de Intercepción No Intrusiva (Efecto Sonda Nulo)}
\label{ssec:efecto_sonda_nulo}

Las técnicas tradicionales de trazabilidad de software basadas en instrumentación de código o puntos de interrupción (\textit{breakpoints}) introducen retardos artificiales que alteran la temporalidad original del programa, un fenómeno conocido como **Efecto Sonda** (\textit{Probe Effect}). En sistemas de tiempo real estricto (\textit{Hard Real-Time}), esta sobrecarga invalida los resultados de la certificación.

Para lograr una intrusión nula, el sistema HR-RT-Reporter implementa una lógica de monitoreo pasivo en silicio que opera como un analizador de bus (\textit{Hardware Sniffer}). La jerarquía de interconexión observa de manera continua y transparente el bus de datos del procesador RISC-V (\texttt{data\_*}). El protocolo de captura se rige por las siguientes etapas arquitectónicas:
\begin{enumerate}
    \item \textbf{Decodificación de Dirección de Frontera:} Dentro del firmware en C ejecutado por el Ibex, cada vez que el algoritmo alcanza un punto de control crítico o un evento de interés, el compilador emite una instrucción estándar de almacenamiento (\textit{Store Word}, \texttt{sw}) apuntando a la dirección de memoria virtual \texttt{0x80000000}.
    \item \textbf{Captura Instantánea en Silicio:} Cuando la lógica de arbitraje del subsistema detecta una validación en el bus con la dirección de frontera (\texttt{data\_req = 1}, \texttt{data\_we = 1} y \texttt{data\_addr = 30'h80000000}), desvía la petición para no alterar el contenido de la BRAM local e impulsa una señal de disparo hacia el módulo Reporter.
    \item \textbf{Registro Asíncrono:} En el flanco de subida del reloj inmediato al disparo (un tiempo transitorio de exactamente 1 ciclo de reloj, $20\text{ ns}$), el \texttt{axi\_reporter} captura y almacena en su registro interno el valor exacto del contador de $64\text{ bits}$ en ese instante, junto con el dato o identificador de evento emitido por el procesador en el bus de escritura (\texttt{data\_wdata}).
\end{enumerate}

Debido a que el procesador RISC-V ejecuta la instrucción de almacenamiento como parte de su flujo natural de memoria de un ciclo sin experimentar ciclos de espera (\textit{wait states}) ni bloqueo de bus, **el overhead de instrumentación en el tiempo de ejecución del algoritmo es estrictamente constante y matemáticamente acotado a cero ciclos de penalización por arbitraje**, garantizando el cumplimiento de la propiedad de Efecto Sonda Nulo.

% =======================================================================
% SECCIÓN 5: Interconexión y Mapeo del Sistema
% =======================================================================
\section{Interconexión AXI y Mapa de Memoria}
\label{sec:interconexion_axi}

Para articular el intercambio de información entre el dominio del maestro general en la Lógica de Procesamiento (PS) y el hardware de verificación instanciado en el tejido de la Lógica Programable (PL), se implementó una red de interconexión jerárquica AMBA AXI4.

\subsection{Configuración de los Bloques AXI}
\label{ssec:config_axi}

La comunicación se centraliza a través del bloque IP \textit{AXI Interconnect}, configurado para expandir el puerto maestro de propósito general del ARM (\texttt{M\_AXI\_GP0}) hacia cuatro puertos esclavos independientes M-AXI4-Lite operando a una frecuencia de bus de $100\text{ MHz}$. Las direcciones físicas de mapeo para los registros de control y datos se estructuraron como se muestra en la tabla \ref{tab:mapa_memoria}.
%de la siguiente manera:
%\begin{itemize}
%    \item \textbf{\texttt{0x40000000} (AXI BRAM Controller):} Proporciona una ventana de acceso de $8\text{ KB}$ que permite a Linux inicializar y leer la memoria local del procesador Ibex.
%    \item \textbf{\texttt{0x41200000} (AXI GPIO 0 - SUT Reset):} Dedicado exclusivamente al control de la señal de reinicio del procesador RISC-V, acoplado a un canal de $1\text{ bit}$ como salida digital de hardware.
%    \item \textbf{\texttt{0x41210000} (AXI GPIO 1 - Debug LEDs):} Canal digital de $4\text{ bits}$ mapeado hacia los diodos emisores de luz físicos de la placa Zybo para diagnóstico visual rápido del estado operativo.
%    \item \textbf{\texttt{0x43C00000} (HR-RT-Reporter IP):} Dirección base del periférico de trazabilidad, a través de cuya estructura de registros de $64\text{ KB}$ el procesador ARM extrae las marcas temporales (\textit{timestamps}) y banderas de estado capturadas por el silicio.
%\end{itemize}

\subsection{Gatillo de Arranque en Hardware (AXI GPIO)}
\label{ssec:gatillo_arranque}

La validez del ensayo temporal depende críticamente de la sincronización en el arranque del Sistema Bajo Prueba (SUT). Si el núcleo RISC-V iniciara su ejecución de manera autónoma al energizar la FPGA, el software crítico correría antes de que el procesador ARM hubiese finalizado el arranque del kernel Linux y preparado el entorno de recolección de datos.

Para otorgar al sistema operativo un control absoluto sobre el inicio del experimento, la línea física de reinicio activo en bajo del SUT (\texttt{rst\_ni}) fue aislada y subordinada al registro de salida del bloque \textit{AXI GPIO 0}. Mediante esta topología, el procesador Ibex es retenido indefinidamente en estado de reinicio durante el encendido del SoC Zynq. 

El inicio del cómputo crítico solo ocurre en el instante en que el ARM ejecuta una escritura explícita en el registro de control de su mapa de memoria (\texttt{0x41200000}), cambiando el nivel lógico del pin a alto (`1`). Este mecanismo actúa como un **"gatillo de arranque por hardware"** no transaccional, asegurando un punto de partida exacto, determinístico y reproducible para la medición del contador temporal.

% =======================================================================
% SECCIÓN 6: Implementación de Software y Despliegue
% =======================================================================
\section{Entorno de Software Operativo}
\label{sec:software_operativo}

La última etapa metodológica comprende el acondicionamiento del entorno operativo Linux en el procesador ARM Cortex-A9 para interactuar con los periféricos sintetizados en la FPGA.

\subsection{Empaquetado de Bootloader y Modificación del Device Tree}
\label{ssec:bootloader_dts}

El despliegue de la arquitectura física se realizó compilando los productos de diseño de Xilinx Vivado (archivo \texttt{.xsa}) mediante las herramientas de Xilinx Vitis. Se generó un archivo binario de arranque maestro (\texttt{BOOT.bin}), integrando el cargador de primera etapa (\textit{First Stage Bootloader} - FSBL), el bitstream de configuración de la FPGA y el gestor de arranque \textit{U-Boot}.

%A nivel del núcleo Linux, se debió abordar una problemática crítica inherente a la administración de energía del sistema operativo. Por defecto, durante la secuencia de arranque, el kernel de Linux escanea el árbol de dispositivos (\textit{Device Tree} - archivo \texttt{.dts}) y, en un esfuerzo por reducir el consumo energético, desactiva automáticamente todas las líneas de reloj del SoC que no estén asociadas a un controlador de software activo. Dado que el procesador Ibex y el Módulo Reporter operan en la PL sin un driver de kernel monolítico tradicional, Linux apagaría el reloj base de $100\text{ MHz}$, congelando el hardware bajo prueba.

%Para neutralizar este comportamiento, se modificó el archivo de configuración del árbol de dispositivos introduciendo el argumento de arranque \texttt{clk\_ignore\_unused} dentro de las etiquetas \texttt{bootargs}:

%\begin{verbatim}
%/dts-v1/;
%/ {
%    chosen {
%        bootargs = "console=ttyPS0,115200 root=/dev/mmcblk0p2 rw earlycon clk_ignore_unused";
%    };
%};
%\end{verbatim}

%Esta directiva de bajo nivel fuerza al kernel del sistema operativo a mantener energizadas y activas todas las líneas de reloj sintetizadas hacia el tejido programable, garantizando la estabilidad del oscilador de la PL durante la totalidad del experimento.

\subsection{Control AXI desde Espacio de Usuario (User-Space)}
\label{ssec:user_space_io}

En el diseño de sistemas embebidos bajo Linux, la interacción con hardware personalizado puede resolverse mediante la programación de módulos y controladores de espacio de núcleo (\textit{Kernel-space drivers}). Sin embargo, el desarrollo de estos controladores introduce una alta complejidad arquitectónica y sobrecarga en el mantenimiento del sistema operativo.

Dado que la plataforma HR-RT-Reporter delega toda la criticidad temporal y la captura de alta velocidad directamente a la Lógica Programable (PL), las tareas que debe ejecutar el procesador ARM son estructuralmente asíncronas y de baja frecuencia: iniciar el reset y extraer bloques de resultados finalizados. Por consiguiente, se optó por implementar toda la capa de control de software en el **Espacio de Usuario** (\textit{User-Space I/O}).

Esta técnica se implementó valiéndose de la función de llamada al sistema \texttt{mmap}, la cual proyecta un rango de direcciones físicas del bus AXI (expuesto a través de la partición virtual de memoria del kernel \texttt{/dev/mem}) directamente en el espacio de memoria virtual de un proceso de usuario. De esta forma, aplicaciones estándar de alto nivel pueden acceder a los registros físicos de la FPGA mediante punteros de puntero de lectura y escritura directos, logrando una velocidad de transferencia cercana a la nativa del hardware sin incurrir en cambios de contexto hacia el modo supervisor del kernel.

\subsection{Scripts de Validación y Recolección de Datos}
\label{ssec:scripts_recopilacion}

El control operativo del flujo de verificación se automatizó mediante el desarrollo de un conjunto de rutinas y scripts estructurados en el lenguaje de programación Python, los cuales operan en el sistema operativo Linux empotrado. El ciclo procedimental de validación se divide en dos fases ejecutadas por la suite de software:

\begin{enumerate}
    \item \textbf{Fase de Disparo (\texttt{gatillo.py}):} El script principal establece un mapeo de memoria sobre la dirección física \texttt{0x41200000} (AXI GPIO 0) y verifica en primera instancia que la memoria del SUT sea accesible. Una vez preparado el entorno, el script escribe la palabra hexadecimal \texttt{0x00000001} en el bus, elevando la línea \texttt{rst\_ni} y liberando instantáneamente al procesador Ibex para ejecutar el código en su BRAM.
    El fin de este script fue detectar que los buses de comunicación esten llegando al soft-core Ibex, y que la palabra hexadecimal despierte al soft-core.

    \item \textbf{Exploración Básica de Enlaces (\texttt{test\_f1.py}):} 
    Constituye la primera etapa de validación a nivel de software. Es una rutina de bajo nivel diseñada para comprobar la accesibilidad física de los periféricos mapeados. Emplea la librería \texttt{mmap} para leer los primeros bytes de las direcciones de la BRAM (\texttt{0x40000000}) y del GPIO, confirmando que los puentes AXI del SoC Zynq se encuentran operativos y respondiendo tras la configuración del silicio.
    
    \item \textbf{Fase de Recolección (\texttt{debug\_axi.py}):} Una vez que la prueba ha transcurrido, esta rutina realiza lecturas secuenciales sobre la dirección base del periférico \texttt{0x43C00000} (AXI Reporter). El script extrae la parte baja y alta del contador de $64\text{ bits}$ correspondientes a cada evento interceptado, junto con el identificador de registro, y los procesa en espacio de usuario para generar los reportes de trazabilidad de tiempo real exportables a formatos estándar de análisis técnico.

    \item \textbf{Diagnóstico de Arbitraje y Exclusión (\texttt{debug\_axi.py}):} 
    Herramienta enfocada estrictamente en probar la eficacia del multiplexor de seguridad por hardware. El script manipula la señal de reinicio para congelar al núcleo Ibex e inyecta una palabra de prueba en la memoria BRAM (e.g., \texttt{0xDEADBEEF}), asumiendo el control exclusivo del bus. Esta prueba empírica demuestra que el ARM posee privilegios de lectura/escritura únicamente cuando el SUT está detenido, validando así la protección contra colisiones de datos.

    \item \textbf{Validación Integral y Trazabilidad (\texttt{test\_bus.py}):} 
    Representa el ensayo más robusto y automatizado de la arquitectura. Este script mapea simultáneamente la memoria BRAM, el controlador GPIO y el Módulo Reporter. Su flujo operativo permite inyectar un \textit{firmware} de prueba en lenguaje ensamblador RISC-V directamente en el silicio, efectuar el disparo por hardware del Ibex, y realizar lecturas secuenciales sobre la dirección base del periférico \texttt{0x43C00000} (AXI Reporter). Finalmente, el script extrae la parte baja y alta del contador de $64\text{ bits}$ correspondientes a cada evento interceptado y los procesa en el espacio de usuario, generando reportes de trazabilidad en tiempo real listos para su exportación y análisis técnico.

\end{enumerate}

\subsubsection{Conclusiones y Limitaciones de la Arquitectura}
\label{sssec:conclusion_arquitectura}

El análisis derivado de la ejecución del ecosistema de scripts permitió validar, en un entorno de pruebas real, el comportamiento integral de la arquitectura HR-RT-Reporter. Se comprobó exitosamente la interoperabilidad entre el Sistema de Procesamiento (PS) y la Lógica Programable (PL) a través de la red de interconexión AXI, así como la efectividad del multiplexor de exclusión mutua en hardware. Sin embargo, la etapa de integración también expuso las siguientes limitaciones inherentes al diseño de la plataforma experimental:

\begin{enumerate}
    \item \textbf{Capacidad de Almacenamiento Restringida (BRAM):} Para garantizar un determinismo estricto temporal y una baja latencia (accesos en un solo ciclo de reloj de $20\text{ ns}$), el Sistema Bajo Prueba (SUT) opera exclusivamente sobre la memoria interna de la FPGA. Al estar la BRAM dimensionada en $8\text{ KiB}$ ($2048$ celdas de $32\text{ bits}$), existe una limitación rígida sobre la huella de memoria (\textit{memory footprint}) de los algoritmos de control que pueden ser ensayados. El núcleo Ibex carece de acceso a la memoria externa DDR3 para preservar su aislamiento físico.
    
    \item \textbf{Dependencia de Instrumentación Explícita:} Aunque la arquitectura logra mitigar exitosamente el ``efecto sonda`` al eliminar las penalizaciones por bloqueos de bus o interrupciones del sistema operativo, el mecanismo de captura de eventos sigue requiriendo que el código fuente (C/C++) sea instrumentado. El programa debe ejecutar explícitamente una instrucción de almacenamiento (\texttt{sw}) hacia la dirección virtual \texttt{0x80000000}. Esto implica que el binario ensayado difiere ligeramente del binario final de producción, lo cual es una concesión arquitectónica.

    \end{enumerate}


    A continuacion el diagrama en bloques:
    \begin{figure}[htbp]
    \centering
    \includegraphics[width=\textwidth*1,3]{figs/design_3_recortado.pdf}
    
    \caption{Diagrama de bloques general (\textit{Block Design}) sintetizado en AMD Vivado para la arquitectura heterogénea HR-RT-Reporter.}
    \label{fig:vivado_block_design}
\end{figure}
La Figura \ref{fig:vivado_block_design} ilustra la topología completa del diseño por bloques (\textit{Block Design}) implementado en el entorno AMD Vivado. En ella se aprecia la integración física entre el Sistema de Procesamiento (\textit{Processing System} - PS) y los periféricos instanciados en el tejido de la Lógica Programable (\textit{Programmable Logic} - PL).

A nivel microarquitectónico, el esquema se articula a través de los siguientes bloques e interfaces funcionales fundamentales:

\begin{itemize}
    \item \textbf{Procesador Maestro (\texttt{processing\_system7\_0}):} Representa el bloque duro Zynq-7000 PS. Exporta la interfaz maestra de propósito general AXI4-Lite (\texttt{M\_AXI\_GP0}), desde la cual el sistema operativo Linux gobierna la totalidad de los periféricos esclavos en la PL.
    
    \item \textbf{Red de Interconexión AXI (\texttt{ps7\_0\_axi\_periph}):} Actúa como la matriz de conmutación y arbitraje (\textit{AXI Interconnect}). Recibe las peticiones del procesador ARM y las rutea hacia las direcciones físicas de los esclavos mapeados (\texttt{axi\_bram\_ctrl\_0}, \texttt{axi\_gpio\_0}, \texttt{axi\_gpio\_1} y \texttt{axi\_reporter\_0}).

    \item \textbf{Controlador de Memoria (\texttt{axi\_bram\_ctrl\_0}):} Mapeado en la dirección \texttt{0x40000000}, proporciona la interfaz de conversión entre el protocolo AXI4-Lite del ARM y los buses nativos de la memoria BRAM del subsistema Ibex.

    \item \textbf{Gatillo de Reset y Diagnóstico (\texttt{axi\_gpio\_0} y \texttt{axi\_gpio\_1}):}
    \begin{itemize}
        \item \texttt{axi\_gpio\_0} (\texttt{0x41200000}): Canal AXI GPIO de 1 bit cuya salida controla directamente la línea física de reinicio activo en bajo (\texttt{rst\_ni}) del subsistema RISC-V, funcionando como el disparador del experimento controlado desde Linux.
        \item \texttt{axi\_gpio\_1} (\texttt{0x41210000}): Canal de 4 bits acoplado a los diodos emisores de luz (LEDs) de la placa Zybo para la verificación visual del estado operativo del hardware.
    \end{itemize}

    \item \textbf{Subsistema RISC-V Encapsulado (\texttt{ibex\_bram\_wrapper\_bd\_0}):} Módulo envoltorio que aloja al procesador soft-core Ibex, la memoria BRAM de $8\text{ KiB}$ en modo \textit{True Dual Port} y el multiplexor de exclusión mutua en hardware. Este bloque aísla la complejidad interna del procesador RISC-V expuesta hacia el lienzo gráfico de Vivado.

    \item \textbf{Módulo de Trazabilidad (\texttt{axi\_reporter\_0}):} Periférico personalizado mapeado en la dirección \texttt{0x43C00000}. Contiene el contador de alta resolución a $50\text{ MHz}$ ($20\text{ ns}$) y la lógica de intercepción pasiva que captura marcas temporales ante accesos a la dirección virtual de frontera \texttt{0x80000000} sin introducir penalizaciones por contención de buses.
\end{itemize}